% !TeX root = ../papers/nc-p3.tex
\section{整な特異五次点曲線の排除}
\label{sec:singular-integral-quintic-exclusion}
\begin{theorem}\label{thm:singular-integral-quintic-exclusion}
四辺の左右が点ファイバーで、点グラフが整な被約五次曲線と交わらない有限部の和となる標準AS窓は、
定理\ref{thm:elliptic-quintic-regularity}の滑らかな楕円五次の二類で尽きる。
曲線の滑らかさもGorenstein性も仮定しない。
\end{theorem}
\begin{proof}[二次曲面とGorensteinの場合]
曲線を $C$ とする。平面五次なら算術種数6で、非退化整五次のCastelnuovo上界2により全位置で平面となり、
隣接平面の積での次数上限3に反する。算術種数0では滑らかであり、
種数1の特異整曲線は唯一の節点か尖点を持ち、定理\ref{thm:reduced-quartic-point-curves}の固定イデアルの証明で次数によらず除かれる。
以後 $p_a(C)=2$ とする。次数5の $L_i$ は完全な四切断を持ち、$h^0(L_i^2)=9$ である。
実際、整曲線上の階数1無捩れ双対化層の次数は2であり、これに $L_i^{-n}$ を掛けた負次数の層には非零切断がない。
各位置で $C$ を含む既約二次曲面 $Q_i$ がただ一つあり、$L_i\ne L_{i+1}$ である
（偏極が等しければ全双一次式がその曲面の対角で消える）。
滑らかな二次曲面なら曲線の型は $(2,3)$ またはその逆である。
二次錐なら解消上の狭義変換の型は $m s+5f$、$m=1,2$ である。ここで $s$ は負切断である。
$m=2$ なら $s$ と長さ1で交わり、像は頂点で滑らかである。実際錐への写像の局所座標は
$(z,zu,zu^2)$ で、$s$ の方程式 $z$ が枝の一様化元になる。従って狭義変換は元の曲線と同型である。

$m=1$ では曲線を二項形式で
\[
 [Fs^2:Fst:Ft^2:G],\qquad \deg F=3,\quad\deg G=5,\quad (F,G)=1
\]
と書ける（母線への射影は同型で、頂点に写る因子が $F=0$ である）。
頂点以外は滑らかで、頂点の重複度は3、接次元は3である。後者は辺の三アフィン座標の線型結合が
極大イデアルの二乗に入れば、その二項二次係数が $F$ の全零因子上で消えてしまうからである。
この点はGorensteinではない。余次元2のGorenstein曲線イデアルは完全交叉である
（極小Hilbert--Burch解消の末尾階数が型に等しい）。接次元3での二つの方程式はいずれも二次以上から始まり、
一般平面切断での局所交数は4以上になるからである。以上によりGorensteinの場合には $m=1$ は混在しない。

Gorensteinとして $H=\omega_C$ と置く。滑らかな二次曲面の次数2の母線束も
$m=2$ の錐の母線束も $H$ に一致する。両者は次数2で二切断を持つため、
Riemann--Rochと次数零の可逆層の切断条件から従う。これが唯一の次数2の鉛筆である。
一般の滑らかな $p$ に対して、定理\ref{thm:genus-two-quintic-ruling}で使った反対縮約の
三次行列式 $S_p$ と階数低下因子 $D_p$ を取る。この場合にも $D_p$ は $C$ の滑らかな部分に台を持つ長さ2の
Cartier因子で、$\mathcal O_C(D_p)=K^{-1}(p+\sigma^{-1}p)$ である。理由を明記する。
定理\ref{thm:opposite-curve-divisor}の二核ベクトルの共通因子は、$p$ を一般に固定すれば
相異なる位置でそれぞれ滑らかな点 $p,\sigma^{-1}p$ の可逆イデアルである。
残る二階数束の行列式とその束の計算は $C$ 上でそのまま成立する。
特異点 $q$ を固定しても、位置 $q$ での端支持はともに3、反対位置の曲線は空間を張る。
従って $p$ が動くとき行列の一般階数が1以下となることはない
（正規化の滑らかな点の位数を使う同定理の二つの三次元関数空間の積の下界5）。
よって一般の $p$ では $D_p$ が全特異点と錐の頂点を避ける。

$S_p$ の $Q_2$ への制限も非零である。前掲の五次の定理と同様、
恒等零なら余因子行列を曲面または錐の解消のCox環で列と行に分ける。
各々の $C$ 上の写像は $L_1(-p)$ と $L_3(-\sigma^{-1}p)$ による平面射影である。
これらは一般の $p$ で二次曲線には入らない。入れば次数2の鉛筆の一意性から
該当する次数4の束が $H^2$ となり、$p$ は高々一つである。
従って滑らかな二次曲面での両因子の型は $(1,1)$、錐の解消では $s+2f$ である
（片側の母線次数が零なら二次までの一変数表示となり、または次数2の射を経由する）。
よって前者の射影を $L_2$ の三切断でも表すことになる。
この三切断は基点自由な $L_1(-p)$ の三切断の共通倍率であり、その倍率は
$C$ 自身で正則な線束の射である。基点自由な三切断の一つが単元の開集合で比を取ればよい。
従って $L_2(-x)=L_1(-p)$ となる長さ1のCartier因子 $x$ があり、これは滑らかな点である。
この等式が動く $p$ で成り立てば、二つの表示を比較した相異なる二次因子はどちらも唯一の $|H|$ のファイバーになる。
一方の点を固定すると相手も決まるので不可能である（$L_1\ne L_2$）。

従って三次曲面の切断から $C$ を除いて得る残余は前掲定理の通り一本の母線、または解消上の $s+f$ である。
$C$ への制限は $P_2=L_2/H$ の次数3の有効因子 $E_p$ となる。
実際滑らかな二次曲面で残る母線の束は $L_2/H$、錐でも $s+f=(s+2f)-f$ である。
二次小行列が生成する階数低下のイデアルには $S_p$ の微分が入る。
$D_p$ は曲面と $C$ の滑らかな部分にあるため、局所的に $S_p=c e$ として $E_p$ は $D_p$ を含む。
差は滑らかな一点 $r_p$（長さ1の有効Cartier因子）で、
\[
                p+\tau p+r_p\in |P_2 K|,\qquad \tau=\sigma^{-1}.
\]
これは次元1の線型系である。次数3なのでRiemann--Rochで切断はちょうど二つとなる。

この鉛筆を正規化に引き戻して基点を除いた非定数写像を $h$ とする。
一般の $p$ はただ一つのファイバーに入り、その同じファイバーに $\tau p$ も入る。
$\tau=1$ なら一般点で重複度2になり標数零の分離性に反する。
従って $\tau$ の位数は $\deg h\le3$ を割り2か3である。
位数3なら基点はなく、$r_p=\tau^2p$ である。同じ議論を各位置の一般の $p$ に適用すれば
三点の因子が共通なので全ての $P_i K$ が同じになり、$L_i\ne L_{i+1}$ に反する。
位数2なら写像は次数2で、正規化上の固定部は長さ1である。
この基点は $C$ の滑らかな点にある。実際、一般の $p$ で $p,\tau p$ は滑らかであり、
その二点を除いた有効Cartier因子 $r_p$ も滑らかである。従って $C$ 上の鉛筆の固定部は
滑らかな一点で、これを除いた $C\to\mathbb P^1$ は基点自由な次数2の有限平坦射である。
$\tau$ はその被覆対合である。しかし二重被覆の固定因子の長さは6である。
実際、trace分解は $\mathcal O\oplus N^{-1}$ で、$\chi(\mathcal O_C)=-1$ より $\deg N=3$、
固定因子の局所式は反不変生成元 $y=0$、$y^2$ は $N^2$ の切断である。
他方、定理\ref{thm:fixed-differential}の第三項により固定因子は次数 $2p_a-2=2$ となり矛盾する。
以上は特異な二重被覆でも固定イデアルの長さをそのまま数えている。

\smallskip
\noindent\textit{三重の頂点を持つ場合。}
非Gorensteinの場合は全位置で $m=1$ の錐表示になる。唯一の特異点は四移動が固定し、
その接次元は3である。従って命題\ref{prop:ore-mark-recognition}が適用できる。
水平三文字を全位置で同一視して、辺は一定の三文字商関係 $S$ と
$X_i h-h X_{i+1}=d_i(h)$ である。AS窓なら三文字部分も標準ASである。
頂点からの水平射影は全位置で正規化を滑らかな平面二次曲線に同型に写す。

三文字の点グラフが平面全体の非恒等射影変換なら、定理\ref{thm:plane-ore-unmarked}の表示で
点グラフは曲面を含む（零補正は $X=0$ の平面、固定直線上の非零補正は同定理の二次曲面）。
平面上で恒等なら、任意の持ち上がる一般の水平点で前後の高さに同じ量を加えられ、
二次曲線全体が持ち上がる以上やはり曲面を得る。
残る三文字グラフは平面三次上の自己同型である。水平二次曲線を含むので二次曲線と直線の和であり、
二次曲線は唯一だから全位置で同じである。その上の自己同型を $g$ とする。
$g=1$ も高さを同時に動かせるため除かれる。

一般の正規化点 $p$ で二次形式の座標列を $v(p)$ とすると、混合式は高さ $a,b$ に対して
\[
               a\,v(gp)-b\,v(p)=d_i(v(p),v(gp))
\]
となる（座標の持上げを固定する）。$a$ は本来水平束 $\mathcal O_{\mathbb P^1}(2)$ の有理切断である。
Cramerの公式により、その極因子の次数は $v(p)\wedge v(gp)$ の三成分の共通零因子の次数以下である。
局所的に各座標束を自明化して単元となる二次小行列を共通因子除去後に使えばよい。
二次埋込みは閉埋込みなので、この共通零イデアルは $g$ の固定イデアルである。
非恒等な射影直線自己同型の固定長は重複込みで2である。ところが五次の表示で高さは $G/F$、
その極因子の次数は3で、水平切断の加算でも変化しない。これで最後の可能性も消える。
\end{proof}
